\section*{Chapitre \ref{ch:rs:build-a-research-workflow} --- Construire un flux de travail de recherche}

\begin{solution}{exo:rs:build-a-research-workflow:1}
Les mêmes paramètres exécutés avec un code différent donnent des résultats différents~; une clé sans le hachage du code servirait l'ancien résultat après la modification du code (la
fiche de synthèse qui avait gagné un \omterm{def:rs:performance-measurement:hit}{taux de réussite} (hit rate) aurait été lue dans le cache sans lui).
\end{solution}

\begin{solution}{exo:rs:build-a-research-workflow:2}
Afin qu'un changement en amont qui laisse une sortie intermédiaire inchangée s'arrête là~: les clés aval ne dépendent que du contenu des entrées, et non de la manière
dont elles ont été produites (arrêt précoce).
\end{solution}

\begin{solution}{exo:rs:build-a-research-workflow:3}
Instantané, univers, caractéristiques, \omterm{def:rs:anatomy-of-a-predictor:card}{fiches de prédicteur} et niveau~1 sont recalculés, puisque chacun lit l'instantané~; seule la sortie de l'instantané change, si bien que le mélange, le portefeuille et la fiche de synthèse
sont lus depuis le cache.
\end{solution}

\begin{solution}{exo:rs:build-a-research-workflow:4}
Elle impose de tenir la trace de la manière dont chaque résultat a été produit (1), d'enregistrer les résultats intermédiaires (5), de noter les graines (6, en signalant
dans la reproduction les étapes sans graine), et une partie de l'archivage des versions (3, l'enregistrement de l'environnement). Éviter les étapes manuelles (2), le contrôle de version (4), conserver les données des graphiques (7), une structure hiérarchique
des sorties (8), relier les énoncés aux résultats (9) et l'accès public (10) restent des pratiques, même si le manifeste facilite la pratique (9).
\end{solution}

\begin{solution}{exo:rs:build-a-research-workflow:5}
Une modification de l'auxiliaire change la sortie de l'étape sans changer sa clé, si bien que le cache sert un résultat périmé. Hacher les fonctions que l'étape appelle (le code source de son
module, ou la fermeture transitive des fonctions qu'elle utilise), ou versionner explicitement les auxiliaires en les passant comme paramètres.
\end{solution}

\begin{solution}{exo:rs:build-a-research-workflow:6}
L'environnement~: versions des bibliothèques (la sommation de NumPy et les compilations BLAS modifient l'ordre des opérations en virgule flottante), nombre de threads, plateforme. Le
relevé d'environnement du manifeste indique ce qu'il faut comparer~; une reproduction bit à bit entre environnements peut exiger des versions figées et des réductions monothread (livre~4,
chapitre~25).
\end{solution}

\begin{solution}{exo:rs:build-a-research-workflow:7}
Sans graine, \texttt{reproduce} ne renvoie que la fiche de synthèse. Une graine dérivée du hachage des entrées et des paramètres de l'étape est fixée pour un calcul donné~;
elle se reproduit, et elle n'est choisie par personne, si bien qu'elle ne peut pas être ajustée~: un résultat qui dépend de la graine reste mis en évidence en faisant varier les entrées.
\end{solution}

\begin{solution}{exo:rs:build-a-research-workflow:8}
Le contrôle de version conserve le code, et non les données qu'il a lues (le fournisseur les a révisées), ni l'environnement, ni l'ordre d'exécution des cellules, ni les graines,
ni la version ayant produit les chiffres publiés. Un manifeste enregistre toutes ces informations.
\end{solution}

\begin{solution}{pb:rs:build-a-research-workflow:1}
\begin{enumerate}
\item Instantané~; univers (instantané)~; caractéristiques (instantané, univers)~; \omterm{def:rs:anatomy-of-a-predictor:card}{fiches de prédicteur} (instantané, caractéristiques)~; mélange (caractéristiques, univers)~; portefeuille (mélange, univers)~; niveau~1
(portefeuille, instantané)~; fiche de synthèse (niveau~1).
\item La clé~: nom, hachage du code, paramètres, graine, hachages des sorties des entrées. Le manifeste~: pour chaque étape, sa clé, le hachage du code, les paramètres,
la graine, les dépendances, les hachages des entrées, le hachage de la sortie, si l'étape a été exécutée, et l'environnement.
\item Il lit toutes les étapes depuis le cache~: toutes les clés sont stockées.
\item \omterm{def:rs:anatomy-of-a-predictor:ic}{Coefficient d'information de rang} (rank IC) du momentum 0.0076 ($t = 1.39$), celui du retournement 0.0379 ($t = 16.05$) sur 2\,267 jours~; ratio de Sharpe du portefeuille mensuel après coûts 0.08 (erreur type
0.33, intervalle bootstrap $-0.67$ à 0.74), rotation de 3.5~\% par jour.
\item Instantané, univers, caractéristiques, \omterm{def:rs:anatomy-of-a-predictor:card}{fiches de prédicteur} et niveau~1 relancés~; seule la sortie de l'instantané change.
\item Les huit étapes sont relancées, et les huit sorties changent.
\item Mélange, portefeuille, niveau~1 et fiche de synthèse.
\item La fiche de synthèse seule.
\item Le même hachage de sortie pour les huit étapes~: une liste de différences vide.
\item La reproduction signale la fiche de synthèse, et elle seule.
\item Versions de Python et de NumPy, et plateforme~; lorsqu'une reproduction échoue sur une autre machine, ou que des résultats en virgule flottante diffèrent d'une version de bibliothèque à l'autre.
\item \textbf{Résultat nommé.} Une révision hors de l'univers relance cinq étapes et modifie une sortie (celle de l'instantané)~; une révision à l'intérieur relance les huit étapes et modifie toutes leurs
sorties~; reproduire le manifeste depuis zéro renvoie des hachages de sortie identiques pour les huit étapes.
\item Comme un essai de sa famille, avec le hachage des clés d'étapes du manifeste comme identité du code, le hachage de sortie de l'instantané comme identité des données, et ses
métriques.
\item La reproduction, les \omterm{def:rs:point-in-time-data-and-the-biases:lookahead}{données point-in-time} (en partie~: l'identité de l'instantané et ses révisions), le \omterm{def:rs:the-research-process:log}{nombre d'essais} (par l'enregistrement)~; ni le \omterm{def:rs:vectorised-backtests:fidelity}{niveau de fidélité},
ni les coûts à la taille visée, ni si le texte correspond aux chiffres.
\item Le portefeuille est reconstruit chaque mois et le signal de retournement s'estompe en un jour (inadéquation d'horizon du chapitre~6).
\item Sous forme de manifeste avec son cache et son environnement, enregistré dans le \omterm{def:rs:the-research-process:log}{journal de recherche}~: une seule commande pour le relancer, et une comparaison qui indique ce qu'a modifié la révision du fournisseur.
\item Des hachages de code qui suivent les fonctions appelées, un cache partagé, des \omterm{def:rs:build-a-research-workflow:repro}{verrouillages de l'environnement} appliqués, et une validation de revue consignée avec le manifeste.
\item Un résultat que son manifeste reproduit à l'identique, bit à bit, à partir des mêmes données, du même code, des mêmes paramètres, des mêmes graines et du même environnement.
\end{enumerate}
\end{solution}

\begin{solution}{iq:rs:build-a-research-workflow:1}
Données révisées ou faisant l'objet d'un nouvel instantané (\omterm{def:rs:point-in-time-data-and-the-biases:vintage}{retraitements} du fournisseur, corrections du \omterm{def:rs:point-in-time-data-and-the-biases:survivorship}{biais de survie}), code ou versions de bibliothèques modifiés, paramètres différents, tirages aléatoires sans graine, un
univers ou calendrier différent, \omterm{def:rs:point-in-time-data-and-the-biases:lookahead}{biais d'anticipation} corrigé depuis, ou un environnement différent (threads, plateforme).
\end{solution}

\begin{solution}{iq:rs:build-a-research-workflow:2}
Tenir la trace de la manière dont chaque résultat a été produit~; automatiser chaque étape~; versionner le code et les données (instantanés)~; enregistrer les résultats intermédiaires~; fixer une graine pour chaque source d'aléa~; verrouiller
les environnements~; relier chaque chiffre d'un écrit à sa source (les dix règles de Sandve et de ses coauteurs).
\end{solution}

\begin{solution}{iq:rs:build-a-research-workflow:3}
\omterm{def:rs:build-a-research-workflow:repro}{Stockage adressé par contenu}~: indexer chaque étape par le hachage de son code, de ses paramètres, de sa graine et des hachages de contenu de ses entrées~; stocker les sorties sous leurs clés~; ne jamais invalider
explicitement, puisqu'une entrée modifiée produit une nouvelle clé~; construire les clés à partir des hachages de sortie pour l'arrêt précoce~; ramasser les entrées non référencées.
\end{solution}

\begin{solution}{iq:rs:build-a-research-workflow:4}
Reproduction à partir d'un manifeste~; \omterm{def:rs:point-in-time-data-and-the-biases:lookahead}{données point-in-time} et univers~; \omterm{def:rs:the-research-process:log}{nombre d'essais}, ratio de Sharpe déflaté et PBO~; \omterm{def:rs:vectorised-backtests:fidelity}{niveau de fidélité} du \omterm{def:rs:vectorised-backtests:backtest}{backtest}~; coûts et
\omterm{def:rs:capacity-decay-and-crowding:capacity}{capacité} à la taille visée~; stabilité sur les sous-périodes~; concordance entre les chiffres du texte et les sorties.
\end{solution}

\begin{solution}{iq:rs:build-a-research-workflow:5}
Fixer explicitement la graine de chaque générateur (jamais à partir de l'horloge), utiliser un générateur à flux documenté (et non une valeur par défaut de bibliothèque susceptible de changer), consigner la graine et la
version de la bibliothèque, et rendre l'ordre des tirages indépendant des threads.
\end{solution}

\begin{solution}{iq:rs:build-a-research-workflow:6}
Repérer chaque manifeste dont l'étape d'instantané a utilisé les données du fournisseur (l'identité des données du \omterm{def:rs:the-research-process:log}{journal de recherche})~; relancer chacun avec l'instantané retraité~; les comparaisons indiquent quelles étapes
ont vu leurs sorties changer, et l'arrêt précoce indique quels résultats restent intacts.
\end{solution}
